\section{Related Work}
\textbf{Global alignment beyond strings.} The problem we address (\S\ref{sec:problem}) is a global alignment: every allele of one graph is either matched to an allele of the other or reported as a gap. Its sequence analogue is Needleman-Wunsch \cite{NEEDLEMAN1970443}, an $O(mn)$ dynamic program over two strings. Partial-order alignment (POA) \cite{lee2002poa} generalizes this dynamic program to directed acyclic graphs and can align two partial-order graphs to each other, but its DP ranges over the product of the two node sets and targets gene-scale multiple alignments, not chromosome-scale pangenomes with $10^7$ nodes. Exactly aligning two alignments is already NP-complete \cite{kececioglu2004aligning}, and even the asymmetric sequence-to-graph problem becomes NP-complete once edits are allowed on the graph \cite{jain2020complexity}. Read mappers, whether to a sequence \cite{li2018minimap2} or to a graph \cite{garrison2018vg,xu2025parsga}, keep one input a string; they enter this paper only through the linearize-and-map workaround we compare against in \S\ref{sec:eval}, which expands a graph back into its haplotypes and maps each one separately.

\textbf{General graph comparison.} Graph edit distance and biological network alignment compare arbitrary graphs. Exact graph edit distance is NP-hard; its standard approximation reduces it to a bipartite assignment between local substructures solved by Hungarian-type algorithms \cite{riesen2009approximate}, the same reduction PanG2G applies inside each snarl. Network aligners such as IsoRank \cite{singh2008isorank} match protein-interaction networks by combined topological and sequence similarity. Neither exploits the bubble structure of variation graphs or reports base-pair alleles. PanG2G confines the assignment to paired snarls, where it is small and solved exactly.

\textbf{Pangenome graph comparison.} Pangenome graphs are built by vg \cite{garrison2018vg}, minigraph \cite{li2020minigraph}, Minigraph-Cactus \cite{hickey2023pangenome} and PGGB \cite{garrison2023pangenome}. Comparing two of them today goes either through a reference, where \texttt{vg deconstruct} projects each graph's bubbles onto a reference path as a VCF, as in the HPRC release \cite{Liao2023-short}, and the two VCFs are then compared, or through coarse summaries. ODGI \cite{guarracino2022odgi} computes path-level similarity within a single graph; PanGraph \cite{noll2023pangraph} and SplitMEM \cite{marcus2014splitmem} relate pangenomes at gene-block or suffix-structure granularity, not at the allele level. The Pangenome Graph Variation Format (PGVF) \cite{pgvfspec}, a GFA fork, has been proposed to represent graph-to-graph alignments, but no tool computes them. PanG2G is, to our knowledge, the first system that does so at allele resolution.
